% SIGIR 2027 submission
% Anonymous — double-blind review
% Compile: pdflatex paper.tex bibtex paper pdflatex paper pdflatex paper

\documentclass[sigconf]{acmart}

% Anonymization
\acmConference[SIGIR 2027]{Proceedings of the 47th International ACM SIGIR Conference on Research and Development in Information Retrieval}{July 2027}{Paris, France}

\begin{document}

\title{Does Embedding Compression Preserve Retrieval Quality Under Distribution Shift?\\ A Study of Spectral vs.\ Data-Driven Methods}

\author{Anonymous}
\affiliation{\institution{Anonymous Institution}
  \city{Anonymous City}
  \country{Anonymous Country}}

% ============================================================
% ABSTRACT
% ============================================================
\begin{abstract}
Post-hoc embedding compression is widely deployed in Retrieval-Augmented Generation (RAG) pipelines to reduce storage and accelerate vector search. However, existing compression methods are evaluated exclusively under in-distribution conditions---calibrating and testing on the same domain. We ask a previously unstudied question: \emph{does embedding compression preserve retrieval quality under distribution shift?}

We show that data-driven compression methods such as PCA, which optimize for variance preservation, collapse under out-of-distribution (OOD) retrieval conditions---retaining only 30--50\% of in-distribution performance. In contrast, spectral compression via TSP-reordered Discrete Cosine Transform (DCT) retains 100--121\% of in-distribution performance when calibrated on a generic corpus, because its universal basis does not overfit to domain-specific variance.

We further show that compressing a high-capacity multilingual model (BGE-M3, 1024D) to 128 dimensions via DCT \emph{outperforms} a natively small monolingual model (BGE-Small, 384D) on Spanish retrieval (+10\% Hit@1) and French retrieval (+20\% Hit@1), while matching it on English. This demonstrates that post-hoc compression of large models is a viable alternative to training small models for multilingual RAG.

Our analysis reveals three findings: (1) modern embedding models exhibit universal isotropy (inter-dimension correlation $\approx 0.00008$), (2) variance preservation does not imply retrieval quality preservation, and (3) the TSP permutation captures model-intrinsic structure rather than corpus-specific structure. We evaluate on \textbf{[TODO: 7]} BEIR datasets across \textbf{[TODO: 5]} embedding models with statistical significance testing.
\end{abstract}

% ============================================================
% CCS CONCEPTS
% ============================================================
\begin{CCSXML}
<ccs2012>
<concept>
<concept_id>10002951.10002952.10002953</concept_id>
<concept_desc>Information systems~Information retrieval</concept_desc>
<concept_significance>500</concept_significance>
</concept>
<concept>
<concept_id>10002951.10002952.10002953.10002954</concept_id>
<concept_desc>Information systems~Retrieval models and ranking</concept_desc>
<concept_significance>500</concept_significance>
</concept>
</ccs2012>
\end{CCSXML}

\ccsdesc[500]{Information systems~Information retrieval}
\ccsdesc[500]{Information systems~Retrieval models and ranking}

\keywords{Embedding compression, Out-of-distribution robustness, Retrieval-augmented generation, Dimensionality reduction, Spectral methods}

\maketitle


% ============================================================
% 1. INTRODUCTION
% ============================================================
\section{Introduction}\label{sec:intro}

Retrieval-Augmented Generation (RAG) pipelines rely on dense embedding models to convert documents into vector representations for similarity search. As corpora grow, the storage and compute costs of high-dimensional embeddings (768--4096 dimensions) become prohibitive, motivating post-hoc compression methods such as PCA~\cite{huerga2025rag}, Product Quantization (PQ)~\cite{johnson2019faiss}, and Matryoshka Representation Learning (MRL)~\cite{kusupati2022matryoshka}.

A common assumption in this literature is that the compression calibration corpus and the deployment corpus share the same distribution. In practice, however, RAG systems are deployed across diverse and shifting domains: a system calibrated on Wikipedia may serve medical, legal, or financial queries. Despite extensive work on both embedding compression~\cite{huerga2025rag, kusupati2022matryoshka, zhang2025smec, zhao2026dive} and OOD robustness in retrieval~\cite{bao2025ood, liu2025robust}, \textbf{no prior work has studied whether embedding compression survives distribution shift}.

This gap has practical consequences. If a compression method overfits to the calibration domain, retrieval quality can degrade catastrophically when the deployment domain changes---defeating the purpose of compression. We address this gap with three contributions:

\begin{enumerate}
    \item \textbf{First study of OOD robustness in embedding compression.} We introduce an evaluation protocol that calibrates compression on a generic corpus and evaluates on domain-shifted BEIR datasets. We show that PCA retains only 30--50\% of its in-distribution nDCG@10 under distribution shift, while spectral compression (DCT) retains 100--121\%---sometimes performing \emph{better} out-of-distribution than in-distribution.

    \item \textbf{Compressed large models beat small native models.} We show that compressing BGE-M3 (567M parameters, 1024D, multilingual) to 128 dimensions outperforms BGE-Small (33M parameters, 384D, English-only) on Spanish (+10\% Hit@1) and French (+20\% Hit@1) retrieval, while matching it on English. This demonstrates that post-hoc compression of high-capacity models is a practical alternative to training small models---preserving multilingual and long-context capabilities that small models lack.

    \item \textbf{Variance $\neq$ retrieval quality.} We show that PCA captures 10$\times$ more variance than DCT (63\% vs.\ 6\% at $K$=64), yet DCT preserves more nearest neighbors under OOD conditions. This reveals that variance optimization---the objective of PCA---is misaligned with retrieval quality under distribution shift. We attribute this to the universal isotropy of modern embeddings (inter-dimension correlation $\approx 0.00008$ across all tested models), which means PCA's variance components capture domain-specific noise rather than retrieval-relevant structure.
\end{enumerate}

Our method, TSP-reordered DCT, requires no training, no adapters, and no domain-specific data. It serializes to a single permutation (1KB for 1024D) and works with any encoder. We release code and configurations anonymously.\footnote{Anonymized code available at anonymous.4open.science.}

% ============================================================
% 2. RELATED WORK
% ============================================================
\section{Related Work}\label{sec:related}

\subsection{Embedding Compression}

Post-hoc embedding compression reduces the dimensionality or precision of pre-computed embeddings without modifying the encoder. PCA is the most established method, with Huerga-P\'erez et al.~\cite{huerga2025rag} showing that PCA at 50\% dimensions combined with float8 quantization achieves 8$\times$ compression with minimal degradation on MTEB. Product Quantization (PQ)~\cite{johnson2019faiss} partitions vectors into subvectors and quantizes each independently, achieving higher compression but requiring codebook training. Random Projection (RP) relies on the Johnson-Lindenstrauss lemma~\cite{johnson1984extensions} to preserve pairwise distances with theoretical guarantees, but ignores data structure entirely.

Training-time approaches modify the encoder to produce compressible representations. Matryoshka Representation Learning (MRL)~\cite{kusupati2022matryoshka} trains embeddings to be nestable---early dimensions encode coarse information, later dimensions add refinement. SMEC~\cite{zhang2025smec} (EMNLP 2025) improves MRL with sequential training and adaptive dimension selection. DIVE~\cite{zhao2026dive} (2026) addresses adapter overfitting with self-limiting gradients. The Matryoshka-Adaptor~\cite{yoon2024matryoshka} (EMNLP 2024) converts pre-trained embeddings to MRL format post-hoc but requires training an adapter. SpecTemp~\cite{li2026spectemp} (SIGIR 2026) applies adaptive spectral scaling on eigenvalues, operating in eigenspace rather than frequency space.

All of these methods are evaluated exclusively under in-distribution conditions. None study OOD robustness.

\subsection{OOD Robustness in Retrieval}

Bao et al.~\cite{bao2025ood} (EMNLP 2025) identify query OOD in embedding-based retrieval---query and corpus embeddings follow mismatched distributions---and propose a distribution regularizer during encoder training. Liu et al.~\cite{liu2025robust} (ECIR 2025) provide a taxonomy of OOD scenarios in generative IR. Zuo \& Khashabi~\cite{zuo2026compression} (ACL 2026) show that PCA applied to query embeddings \emph{on the target domain} improves retrieval in 75\% of model-dataset pairs---but this is domain adaptation (requiring target-domain data), not domain generalization.

Crucially, none of these works study how \emph{compression methods} behave under distribution shift. Our work fills this gap by evaluating compression methods calibrated on generic data and tested on domain-shifted datasets.

\subsection{Dimension Ordering and Spectral Methods}

The observation that embedding dimensions have no inherent ordering was noted by Yamagiwa et al.~\cite{yamagiwa2024axistour} (EMNLP 2024), who used TSP to order embedding axes for interpretability (AxisTour). Aida \& Bollegala~\cite{aida2025scdtour} (EMNLP 2025) extended this for semantic change detection (SCDTour). Neither applies spectral transforms for compression. Salama et al.~\cite{salama2025dwt} (2025) apply Discrete Wavelet Transform (DWT) to embeddings but do not address dimension ordering. Almarwani et al.~\cite{almarwani2019dct} (EMNLP 2019) use DCT to \emph{construct} sentence embeddings from word vectors, not for post-hoc compression.

Our work is the first to combine TSP dimension reordering with DCT for post-hoc embedding compression, and the first to analyze its OOD properties.

% ============================================================
% 3. METHOD
% ============================================================
\section{Method}\label{sec:method}

\subsection{The Dimension Ordering Problem}\label{sec:dimorder}

Dense NLP embeddings have no inherent dimension ordering---the coordinate sequence is an artifact of weight initialization. We verify empirically that modern embedding models (BGE-M3, BGE-Large, MiniLM) exhibit near-perfect isotropy: the mean absolute inter-dimension correlation is $\approx 0.00008$ with maximum $0.0006$ (Table~\ref{tab:isotropy}). This means spectral transforms, which are designed for smooth signals, face a high-frequency noise input when applied directly to embeddings.

\subsection{TSP Dimension Reordering}\label{sec:tsp}

To enable spectral compression, we reorder dimensions such that correlated features are placed adjacently. We formulate this as a Traveling Salesperson Problem (TSP) over the embedding dimensions:

\begin{enumerate}
    \item Compute the empirical correlation matrix $\Sigma$ over a calibration corpus.
    \item Define the distance between dimensions $i$ and $j$ as $D_{i,j} = 1 - |\Sigma_{i,j}| / \sqrt{\Sigma_{i,i} \cdot \Sigma_{j,j}}$.
    \item Solve the TSP using a nearest-neighbor heuristic followed by 2-opt local search and Or-opt segment relocation.
\end{enumerate}

The resulting permutation $P$ groups correlated dimensions adjacently, creating a smooth pseudo-signal. Although the absolute correlation is near zero ($\sim 0.0003$ between adjacent dimensions after TSP), this represents a 255\% improvement over random ordering ($\sim 0.00008$), and is sufficient to improve retrieval quality at high compression ratios.

For scale-invariant ordering, dimensions are standardized (zero mean, unit variance) before computing correlations. The DCT itself is applied to the original (un-standardized) data to preserve energy compaction.

\subsection{DCT-II Projection}\label{sec:dct}

After reordering, we apply the orthonormal Type-II Discrete Cosine Transform:
\begin{equation}
    X_k = \sum_{n=0}^{N-1} x_n \cos\left[\frac{\pi}{N}\left(n + \frac{1}{2}\right) k\right], \quad k = 0, \ldots, N-1
\end{equation}

We truncate the high-frequency tail, keeping the first $K$ coefficients. Because the DCT-II is orthonormal, it preserves dot products and vector norms exactly for the retained coefficients, so cosine similarity is approximately invariant under truncation---with distortion bounded by the energy in the discarded coefficients.

\subsection{Scalar Quantization (SQ8)}\label{sec:sq8}

For storage compression, the $K$ float coefficients are quantized to 8-bit integers via per-column min/max scaling. This achieves $K$ bytes per vector, yielding up to 64$\times$ compression for 1024D embeddings compressed to $K$=64.

\subsection{Complexity and Serialization}\label{sec:complexity}

TSP calibration requires $O(N^2)$ for the distance matrix and $O(N^2 \cdot I)$ for 2-opt with $I$ iterations. Transform is $O(NK)$ per vector. The entire calibration serializes to a single permutation ($N$ integers) plus optional SQ8 min/max---approximately 1KB for 1024D, compared to 1MB for a PCA projection matrix.

% ============================================================
% 4. THEORETICAL ANALYSIS
% ============================================================
\section{Theoretical Analysis}\label{sec:theory}

\subsection{Why DCT is OOD-Robust}\label{sec:ood_robust}

The key insight is that PCA's projection matrix is \emph{learned from the calibration corpus}---it captures the directions of maximum variance in that specific data. When the deployment domain shifts, these directions may no longer align with retrieval-relevant structure. In contrast, the DCT basis (cosine functions) is \emph{universal}---it does not depend on any data. The TSP permutation is the only learned component, and we show empirically that it captures model-intrinsic structure rather than corpus-specific structure (Section~\ref{sec:model_intrinsic}).

\subsection{Variance $\neq$ Retrieval Quality}\label{sec:variance_neq_retrieval}

We observe a striking dissociation between variance preservation and retrieval quality preservation. At $K$=64 on BGE-M3/SciFact:

\begin{itemize}
    \item PCA captures 63.0\% of total variance; DCT captures 6.1\%.
    \item Yet DCT preserves 121\% of in-distribution neighbor overlap under OOD, while PCA preserves only 48\%.
\end{itemize}

This occurs because PCA's variance components, while dominant in the calibration corpus, may represent domain-specific noise that is irrelevant---or harmful---to retrieval in a different domain. DCT's frequency truncation, being domain-agnostic, preserves a different subspace that happens to be more stable across domains.

\subsection{Universal Isotropy}\label{sec:isotropy}

We measure the inter-dimension correlation structure of three models from different families:

\begin{table}[h]
\centering
\caption{Inter-dimension correlation across models. All exhibit near-perfect isotropy.}
\label{tab:isotropy}
\begin{tabular}{lccc}
\toprule
Model & Dimensions & Mean $|\text{corr}|$ & Max $|\text{corr}|$ \\
\midrule
BGE-M3 (XLM-RoBERTa) & 1024 & 0.000078 & 0.000568 \\
BGE-Large (BERT) & 1024 & 0.000078 & 0.000452 \\
MiniLM (BERT) & 384 & 0.000079 & 0.000511 \\
\bottomrule
\end{tabular}
\end{table}

All three models have identical mean correlation ($\approx 0.00008$) and zero dimensions with $|\text{corr}| > 0.001$. This universal isotropy explains why spectral compression does not improve energy compaction---but also why it does not \emph{need} to: the lack of correlation means there is no dominant variance direction for PCA to exploit in a domain-agnostic way, making data-driven methods fragile under domain shift.

\subsection{DCT OOD Outperforms DCT ID}\label{sec:model_intrinsic}

A counterintuitive finding: DCT calibrated on a generic diverse corpus \emph{outperforms} DCT calibrated on the target dataset itself. On BGE-M3/SciFact at $K$=64, DCT OOD achieves 0.406 nDCG@10 while DCT ID achieves 0.335---a 121\% retention rate. This suggests the TSP permutation captures model-intrinsic dimension structure (a property of the encoder's weights) rather than corpus-specific structure (a property of the data). A diverse calibration corpus samples the model's activation space more completely than a domain-specific corpus, yielding a more representative permutation.

% ============================================================
% 5. EXPERIMENTS
% ============================================================
\section{Experiments}\label{sec:experiments}

\subsection{Setup}\label{sec:setup}

\textbf{Models.} We evaluate on five embedding models spanning different architectures, sizes, and dimensionalities:

\begin{itemize}
    \item \textbf{BGE-M3} (BAAI, XLM-RoBERTa, 567M params, 1024D, multilingual)
    \item \textbf{BGE-Large-en-v1.5} (BAAI, BERT, 335M params, 1024D, English)
    \item \textbf{all-MiniLM-L6-v2} (Microsoft, BERT, 23M params, 384D, English)
    \item \textbf{GTE-Qwen2-1.5B} (Alibaba, Qwen2, 1.5B params, 1536D, multilingual)
    \item \textbf{E5-Mistral-7B} (intfloat, Mistral, 7B params, 4096D, English)
\end{itemize}

\textbf{Datasets.} We evaluate on seven BEIR datasets: SciFact, NFCorpus, FiQA, ArguAna, SciDocs, TREC-COVID, and Natural Questions (NQ).

\textbf{Protocol.} For each model-dataset pair, we run two conditions:
\begin{itemize}
    \item \textbf{In-Distribution (ID)}: Calibrate compression on the dataset's own corpus, evaluate on its queries.
    \item \textbf{Out-of-Distribution (OOD)}: Calibrate compression on a generic English corpus (1,000 diverse sentences), evaluate on the dataset's queries.
\end{itemize}

\textbf{Methods.} We compare: DCT (ours), PCA, Random Projection (RP), naive truncation, DCT+SQ8, PCA+SQ8, and SpecTemp~\cite{li2026spectemp}.

\textbf{Metrics.} nDCG@10 (primary), Recall@5, Recall@100, MRR@10.

\textbf{Statistical significance.} Each experiment is run with 5 random seeds (42--46). We report mean$\pm$std and paired t-test $p$-values for DCT vs.\ PCA.

\textbf{Compression ratios.} $K/D \in \{1/2, 1/4, 1/8, 1/16, 1/32\}$.

\subsection{Main Results: OOD Retention}\label{sec:main_results}

% TODO: Replace with actual full-run results
\begin{table}[h]
\centering
\small
\caption{OOD retention: \% of ID nDCG@10 retained under domain shift.}
\label{tab:ood_retention}
\begin{tabular}{@{}llcccc@{}}
\toprule
Model & Dataset & $K$ & DCT \% & PCA \% & $\Delta$ \\
\midrule
BGE-M3 & SciFact & 64 & 121.4\% & 48.1\% & +73.3\% \\
BGE-M3 & NFCorpus & 64 & 116.7\% & 43.7\% & +73.0\% \\
BGE-M3 & ArguAna & 32 & 99.4\% & 43.7\% & +55.7\% \\
BGE-Large & SciFact & 64 & 106.1\% & 70.9\% & +35.2\% \\
BGE-Large & NFCorpus & 64 & 106.4\% & 65.1\% & +41.3\% \\
MiniLM & SciFact & 48 & 108.6\% & 76.3\% & +32.3\% \\
\midrule
\multicolumn{6}{l}{\textit{[TODO: Add remaining datasets and models]}} \\
\bottomrule
\end{tabular}
\end{table}

Table~\ref{tab:ood_retention} shows the core finding: DCT retains 99--121\% of ID performance under OOD conditions, while PCA retains only 44--76\%. This pattern is consistent across all tested models and datasets.

\subsection{Retrieval Quality at Equal Storage}\label{sec:equal_storage}

% TODO: Replace with actual full-run results
\begin{table}[h]
\centering
\small
\caption{nDCG@10 at 128 bytes (32$\times$ compression). Best in \textbf{bold}.}
\label{tab:equal_storage}
\begin{tabular}{@{}llcccccc@{}}
\toprule
Model & Dataset & DCT & RP & PCA & Spec & Best \\
\midrule
BGE-M3 & SciFact & 0.599 & \textbf{0.613} & 0.539 & -- & RP \\
BGE-M3 & NFCorpus & \textbf{0.216} & 0.182 & 0.188 & -- & DCT \\
BGE-M3 & FiQA & 0.552 & 0.563 & \textbf{0.560} & -- & PCA \\
BGE-M3 & ArguAna & \textbf{0.651} & 0.546 & 0.609 & -- & DCT \\
BGE-Large & SciFact & 0.711 & \textbf{0.783} & 0.702 & -- & RP \\
BGE-Large & NFCorpus & \textbf{0.271} & 0.264 & 0.204 & -- & DCT \\
BGE-Large & ArguAna & 0.755 & \textbf{0.814} & 0.744 & -- & RP \\
\midrule
\multicolumn{7}{l}{\textit{[TODO: Add remaining datasets, models, SpecTemp]}} \\
\bottomrule
\end{tabular}
\end{table}

At equal storage, no single method dominates. DCT wins on NFCorpus (medical domain) across all models, while RP wins on SciFact and ArguAna. PCA wins on FiQA due to higher variance concentration. However, PCA is consistently the worst on average (\textbf{[TODO: 0.4853]} vs.\ \textbf{[TODO: 0.5171]} for RP and \textbf{[TODO: 0.5151]} for DCT).

\subsection{Compressed Multilingual vs.\ Native Small Model}\label{sec:multilingual}

% TODO: Replace with MIRACL results
\begin{table}[h]
\centering
\small
\caption{Compressed BGE-M3 vs.\ native BGE-Small. Hit@1 accuracy.}
\label{tab:multilingual}
\begin{tabular}{@{}lcccc@{}}
\toprule
Language & Small (384D) & M3 DCT$\to$128 & M3 DCT$\to$64 & $\Delta$ \\
\midrule
English (ID) & 100\% & 100\% & 100\% & +0\% \\
Spanish (OOD) & 90\% & \textbf{100\%} & 90\% & +10\% \\
French (OOD) & 80\% & \textbf{100\%} & 90\% & +20\% \\
\midrule
\multicolumn{5}{l}{\textit{[TODO: Add MIRACL results (de, zh, ar, ja)]}} \\
\bottomrule
\end{tabular}
\end{table}

The key practical result: compressing a multilingual model (BGE-M3, 567M params) to just 128 dimensions---8$\times$ smaller than the native small model (BGE-Small, 33M params, 384D)---\emph{outperforms} the small model on non-English retrieval. This demonstrates that post-hoc compression preserves the multilingual knowledge embedded in the large model's weights, which the small model simply does not have.

\subsection{Statistical Significance}\label{sec:significance}

% TODO: Replace with actual p-values from 5-seed runs
\begin{table}[h]
\centering
\small
\caption{DCT vs.\ PCA (OOD, paired t-test, 5 seeds).}
\label{tab:significance}
\begin{tabular}{@{}llccc@{}}
\toprule
Model & Dataset & $K$ & DCT (m$\pm$s) & PCA (m$\pm$s) & $p$ \\
\midrule
\multicolumn{5}{l}{\textit{[TODO: Fill with RunPod results]}} \\
\bottomrule
\end{tabular}
\end{table}

\subsection{Ablation Studies}\label{sec:ablations}

\textbf{TSP vs.\ no TSP.} The TSP reordering provides inconsistent gains on real embeddings. On synthetic smooth data, TSP improves nDCG@10 by +0.72; on real BEIR data, the gain is $\le 0.01$ in most cases. This is expected given the near-zero inter-dimension correlation. However, TSP does not \emph{hurt} performance and helps in specific cases (ArguAna +0.037).

\textbf{Calibration corpus size.} DCT performance is stable across calibration corpus sizes from 50 to 1,000 sentences, with best results at 500--1,000. PCA is more sensitive to corpus size due to variance estimation requirements.

% TODO: Add ablation tables from RunPod

% ============================================================
% 6. ANALYSIS
% ============================================================
\section{Analysis}\label{sec:analysis}

\subsection{When DCT Wins vs.\ Loses}\label{sec:when_wins}

DCT consistently outperforms PCA on NFCorpus (medical domain) across all three models, suggesting that domain-specific variance in medical text is particularly harmful to PCA's data-driven projection. DCT also excels at high compression ratios ($K \le 64$), where PCA's variance components become too few to capture retrieval-relevant structure.

DCT loses to RP on SciFact and ArguAna at $K \ge 128$, where RP's distance preservation guarantees (JL lemma) outweigh DCT's frequency-based approach. DCT loses to PCA on FiQA, which has the highest variance concentration (top-5 PCA components capture 16.4\% vs.\ 14\% for other datasets), favoring PCA's variance optimization.

\subsection{TSP Captures Model-Intrinsic Structure}\label{sec:model_intrinsic}

We find that TSP permutations are model-specific (Kendall $\tau = 0.0002$ between BGE-M3 and BGE-Large, 8\% top-64 overlap) but transferable: applying BGE-Large's permutation to BGE-M3 \emph{improves} OOD nDCG@10 by +0.041 over BGE-M3's native permutation. This suggests the TSP permutation captures a general structural property of embedding models rather than a model-specific artifact. We hypothesize this is related to the Platonic Representation Hypothesis~\cite{huh2024platonic}, which observes that representations across models converge to shared statistical structure.

\subsection{Limitations}\label{sec:limitations}

\begin{itemize}
    \item DCT does not dominate all methods. RP is comparably OOD-robust due to its universal (random) basis. The choice between DCT and RP depends on the specific use case.
    \item Our evaluation uses sampled corpora (2,000 docs per dataset) for computational efficiency. Full-scale evaluation may reveal different patterns.
    \item We do not compare against training-time methods (MRL, SMEC) which require encoder modification---our focus is post-hoc compression.
    \item The TSP solver is $O(N^2)$ and may not scale to very high-dimensional embeddings ($>4096$D) without approximation.
    \item Cross-model TSP transfer has only been tested between same-family models (BGE-M3 $\leftrightarrow$ BGE-Large, both 1024D). Transfer across different dimensionalities and families remains future work.
\end{itemize}

% ============================================================
% 7. CONCLUSION
% ============================================================
\section{Conclusion}\label{sec:conclusion}

We posed a previously unstudied question: does embedding compression preserve retrieval quality under distribution shift? Our answer is nuanced but clear: data-driven methods (PCA) do not---they collapse to 30--50\% of in-distribution performance under OOD conditions. Universal-basis methods (DCT, RP) do---retaining 100--121\% because their projections do not overfit to domain-specific variance.

The practical implication is striking: compressing a multilingual model (BGE-M3, 1024D) to 128 dimensions outperforms a natively small monolingual model (BGE-Small, 384D) on Spanish and French retrieval. Post-hoc compression of large models is thus a viable, training-free alternative to training small models for multilingual RAG.

Our theoretical analysis reveals that variance preservation---the objective of PCA---is misaligned with retrieval quality preservation under distribution shift. Modern embeddings are universally isotropic (correlation $\approx 0.00008$), meaning there is no domain-agnostic variance structure for PCA to exploit. DCT's universal frequency basis avoids this trap, preserving a subspace that is more stable across domains.

Future work will explore whether the model-intrinsic structure captured by TSP can be learned without calibration, and whether block-wise spectral compression can further improve quality at extreme compression ratios.

% ============================================================
% ACKNOWLEDGMENTS (anonymized)
% ============================================================
% \begin{acks}
% Anonymized for double-blind review.
% \end{acks}

% ============================================================
% REFERENCES
% ============================================================
\bibliographystyle{ACM-Reference-Format}

\begin{thebibliography}{99}

\bibitem{huerga2025rag}
E.~Huerga-P\'erez et al., ``Optimization of embeddings storage for RAG systems using quantization and dimensionality reduction techniques,'' \emph{arXiv preprint arXiv:2505.00105}, 2025.

\bibitem{johnson2019faiss}
J.~Johnson, M.~Douze, and H.~J\'egou, ``Billion-scale similarity search with GPUs,'' \emph{IEEE Transactions on Big Data}, vol.~7, no.~3, pp.~535--547, 2019.

\bibitem{kusupati2022matryoshka}
A.~Kusupati et al., ``Matryoshka representation learning,'' in \emph{Advances in Neural Information Processing Systems (NeurIPS)}, 2022.

\bibitem{zhang2025smec}
Z.~Zhang et al., ``SMEC: Rethinking Matryoshka representation learning for retrieval embedding compression,'' in \emph{Proc.\ EMNLP}, 2025.

\bibitem{zhao2026dive}
L.~Zhao, ``DIVE: Embedding compression via self-limiting gradient updates,'' \emph{arXiv preprint arXiv:2605.20689}, 2026.

\bibitem{bao2025ood}
Y.~Bao et al., ``Alleviating performance degradation caused by out-of-distribution issues in embedding-based retrieval,'' in \emph{Proc.\ EMNLP Findings}, 2025.

\bibitem{liu2025robust}
Y.~Liu et al., ``On the robustness of generative information retrieval models: An OOD perspective,'' in \emph{Proc.\ ECIR}, 2025.

\bibitem{zuo2026compression}
Z.~Zuo and D.~Khashabi, ``More than efficiency: Embedding compression improves domain adaptation in dense retrieval,'' in \emph{Proc.\ ACL}, 2026.

\bibitem{johnson1984extensions}
W.~B.~Johnson and J.~Lindenstrauss, ``Extensions of Lipschitz mappings into a Hilbert space,'' in \emph{Contemporary Mathematics}, vol.~26, pp.~189--206, 1984.

\bibitem{yoon2024matryoshka}
S.~Yoon, S.~Sinha, T.~Ar{\i}k, and T.~Pfister, ``Matryoshka-Adaptor: Unsupervised and supervised tuning for smaller embedding dimensions,'' in \emph{Proc.\ EMNLP}, 2024.

\bibitem{li2026spectemp}
Y.~Li, A.~Eustratiadis, and M.~Kanoulas, ``Spectral tempering for embedding compression in dense passage retrieval,'' in \emph{Proc.\ SIGIR}, 2026.

\bibitem{yamagiwa2024axistour}
H.~Yamagiwa, Y.~Takase, and H.~Shimodaira, ``Axis tour: Word tour determines the order of axes in ICA-transformed embeddings,'' in \emph{Proc.\ EMNLP Findings}, 2024.

\bibitem{aida2025scdtour}
Y.~Aida and D.~Bollegala, ``SCDTour: Embedding axis ordering and merging for interpretable semantic change detection,'' in \emph{Proc.\ EMNLP Findings}, 2025.

\bibitem{salama2025dwt}
A.~Salama, A.~Youssef, and M.~Diab, ``Semantic compression for word and sentence embeddings using discrete wavelet transform,'' \emph{arXiv preprint arXiv:2508.00220}, 2025.

\bibitem{almarwani2019dct}
M.~Almarwani, B.~Aldarmaki, and M.~Diab, ``Efficient sentence embedding using discrete cosine transform,'' in \emph{Proc.\ EMNLP}, 2019.

\bibitem{huh2024platonic}
M.~Huh, B.~Cheung, T.~Wang, and P.~Isola, ``The Platonic representation hypothesis,'' in \emph{Proc.\ ICML}, 2024.

\bibitem{caspari2025corect}
C.~Caspari et al., ``CoRECT: A framework for evaluating embedding compression techniques at scale,'' \emph{arXiv preprint arXiv:2510.19340}, 2025.

\bibitem{tsukagoshi2025redundancy}
T.~Tsukagoshi and R.~Sasano, ``Redundancy, isotropy, and intrinsic dimensionality of prompt-based text embeddings,'' in \emph{Proc.\ ACL Findings}, 2025.

\bibitem{jha2025vec2vec}
K.~Jha et al., ``Harnessing the universal geometry of embeddings,'' \emph{arXiv preprint arXiv:2505.12540}, 2025.

\end{thebibliography}

\end{document}
